Supervised classifiers trained on Monte Carlo configurations from deep inside two phases and then applied across the transition have become a standard demonstration of machine learning for statistical physics. In the 2D Ising model, Carrasquilla and Melko~\cite{carrasquilla2016machine} trained a feed-forward network on Ising configurations labelled as ordered or disordered by the known critical temperature, and located the critical point from the crossing of the network output at several lattice sizes, together with a finite-size-scaling collapse for $T_c$ and the correlation-length exponent. Alternatives need no phase labels at all: principal component analysis (PCA) of the configurations yields a leading component that tracks the magnetisation~\cite{wang2016discovering,hu2017discovering}, and learning by confusion~\cite{nieuwenburg2016learning} sweeps a trial critical point and looks for a W-shaped accuracy curve whose middle peak marks the transition. The exact answer is known (Onsager~\cite{onsager1944crystal}), $T_c=2/\ln(1+\sqrt2)$, so the 2D Ising model is a convenient place to ask a controlled question: \emph{how large is the error of each read-out, on identical data, and what drives it?}

In this study we want to separate several things that a single demonstration leaves entangled: the finite-size shift of a crossing; the choice of classifier, where a linear model on raw spins cannot see the order parameter because the Hamiltonian is symmetric under a global spin flip; the choice of training window; the number of training samples; and the extra step of an FSS collapse. Related studies examine what trained networks learn~\cite{suchsland2018parameter} and when unsupervised methods work~\cite{hu2017discovering}, but we are not aware of a side-by-side comparison of all of these read-outs under one protocol and one metric.

\paragraph*{What we do.} We generate Ising samples with a batched cluster sampler (Sec.~\ref{sec:setup}), reimplement the supervised classifiers, PCA and learning by confusion, add a non-learned susceptibility/Binder-cumulant reference, and measure the absolute error of the estimated critical temperature against the Onsager value for $L=16,24,32$ and after an FSS collapse, over \rhval{const/n_seeds} seeds with data regenerated for each seed. We register six hypotheses before the main runs (Sec.~\ref{sec:method}) and report each verdict, including the failures.

\paragraph*{Findings, in brief.} (i) The nonlinear classifiers give single-size crossings within roughly half a temperature-grid step of $T_c$ at $L=32$, better than PCA and not distinguishable from confusion, which is nominally lowest at $L=24$ and $L=32$. (ii) Logistic regression on raw spins fails for a symmetry reason that our gauge-fixing control isolates. (iii) The FSS collapse does not help here, and the collapse exponent of the CNN is biased. (iv) The MLP error depends on the training-set size more than on the width of the excluded window around $T_c$; for the CNN the apparent dependence is an optimisation-budget effect. This is a small CPU study; Sec.~\ref{sec:limitations} lists what it cannot say.
